% ID: SEL-D0020-OA-E01
% ESTADO: validada. S021, revisada con los criterios oficiales.
Sean $x$ e $y$ los números de collares A y B. Al dividir las
restricciones de perlas por $10$, $20$ y $30$, respectivamente,
obtenemos $2x+y\le90$, $x+y\le70$ y $x+2y\ge60$, con $x,y\ge0$.
En el eje vertical queda $30\le y\le70$ y el eje horizontal
no contiene puntos factibles. Las fronteras $2x+y=90$ y $x+y=70$
se cortan en $(20,50)$, y $2x+y=90$ y $x+2y=60$ en $(40,10)$.
La otra intersección tiene $y<0$. Estos son los cuatro vértices:
\[
\begin{array}{c|rrrr}
(x,y)&(0,30)&(0,70)&(20,50)&(40,10)\\ \hline
600x+500y&15000&35000&37000&29000
\end{array}
\]

\begin{center}
\setlength{\unitlength}{0.95mm}
\begin{picture}(53,78)(-5,-3)
\qbezier(0,0)(24,0)(48,0)\put(49,0){$x$}
\qbezier(0,0)(0,37)(0,74)\put(0,74){$y$}
\qbezier(0,30)(0,50)(0,70)
\qbezier(0,70)(10,60)(20,50)
\qbezier(20,50)(30,30)(40,10)
\qbezier(40,10)(20,20)(0,30)
\put(2,31){$(0,30)$}\put(2,70){$(0,70)$}
\put(21,51){$(20,50)$}\put(40,12){$(40,10)$}
\end{picture}
\end{center}

Una función lineal sobre este polígono acotado alcanza su máximo
en un vértice. El vértice ganador tiene coordenadas enteras, así
que también resuelve el problema de producción:
$\boxed{20\text{ de A},\ 50\text{ de B},\ 37000\text{ euros}}$.

Para $(40,20)$ se requieren $20\cdot40+10\cdot20=1000$ perlas blancas.
La disponibilidad es de 900, de modo que esa producción es inviable.
